\chapter{Background}

\section{Autonomous Drone Systems}

Autonomous drone systems are becoming increasingly prominent in applications that are otherwise time-consuming or dangerous using traditional methods.
Such applications include infrastructure inspection, for example monitoring bridge structures \cite{ellenberg2016bridge}, environmental monitoring such
as wildlife surveys \cite{anderson2013uav}, and search and rescue (SAR) missions, where drones can assist in locating people after disasters \cite{erdelj2016uav}.
Unmanned Aerial Vehicles (UAVs) are also widely used in surveillance and security applications, including border monitoring and crowd observation \cite{zhou2023uav}.
In addition, drones have been successfully deployed for logistics and delivery tasks, such as Amazon Prime Air \cite{amazon_prime_air}, as well as real-world medical
 transport solutions, including long-range delivery of blood samples in Denmark \cite{investindk2022health_drone}.
Large-scale industrial and energy infrastructures can likewise be efficiently monitored using UAVs \cite{nooralishahi2021drone}.

These diverse applications highlight the potential of UAVs; however, understanding the characteristics and capabilities of different drone platforms is essential.

Multirotor drones are widely used due to their ability to hover, perform vertical takeoff and landing (VTOL), and navigate precisely in constrained environments.
These capabilities make them particularly suitable for urban operations such as inspection and delivery.
In contrast, fixed-wing drones are designed for long-endurance missions and can cover large areas efficiently, making them well suited for large-scale mapping and
environmental monitoring tasks. However, they lack hovering capability and require more space for takeoff and landing, limiting their use in high-precision or urban scenarios.

The suitability of a drone for a given mission depends heavily on the payload it carries.
RGB cameras provide visual data for inspection, surveillance, and mapping tasks, and can also be used for motion estimation through visual odometry \cite{rostum2023visual}.
Thermal cameras mounted on UAVs have been successfully applied in infrastructure inspection, for example to detect subsurface delamination in bridge decks using infrared 
thermography \cite{ellenberg2016bridge}.
LiDAR-equipped UAVs enable accurate three-dimensional mapping of environments, for instance in forestry applications where canopy structure can be estimated from aerial point 
clouds \cite{wallace2016assessment}.
In addition, onboard sensors such as Inertial Measurement Units (IMUs), GPS, and barometers are used for state estimation, including position, velocity, and altitude.

In multi-drone systems, payload diversity introduces functional heterogeneity, meaning that not all drones are capable of performing all tasks. 
This has direct implications for task allocation, as tasks must be matched with drones that possess the required sensing and operational capabilities.

Many modern UAV applications rely on artificial intelligence (AI) for tasks such as perception, navigation, and decision-making. 
However, developers often face strict hardware constraints in drone systems. While cloud-based AI can provide significant computational resources, it introduces latency and 
dependence on network availability.
The emergence of edge AI enables real-time decision-making by allowing drones to process sensor data locally and operate independently of external infrastructure
\cite{mcenroe2022edgeaiuav}. These platforms often support hardware acceleration through GPUs, DSPs, or NPUs, enabling efficient execution of machine learning models.
This is particularly relevant for deploying compact and quantized large language models (LLMs) directly onboard drones.

\section{Multi-Agent Systems}

Multi-agent systems (MAS) consist of multiple autonomous agents that interact in a shared environment to achieve a collective goal \cite{wooldridge2002multiagent}. 
In the context of autonomous drone systems each drone can be modelled as an autonomous agent with it's own sensing and decision making capabilities.
Deploying heterogeneous UAV agents increases the number of available data modalities, the area that the drone swarm can inspect and provides backup options in case of
drone failures, thereby increasing mission success rate in constrast to a single agent solution \cite{pantelimon2019survey}.
The key challenges in multi-agent systems are the efficient coordination between agents and the creation of a mission plan for heterogeneous agents,
particularly under constraints such as limited communication, energy, and computation.

\section{Multi-Agent Systems}

Multi-agent systems (MAS) consist of multiple autonomous agents that interact within a shared environment to achieve individual or collective goals. 
In the context of autonomous drone systems, each drone can be modeled as an autonomous agent with its own sensing, computation, communication, and decision-making capabilities. 
Deploying heterogeneous UAV agents can increase the number of available sensing modalities, expand the area that can be inspected, and provide redundancy in case of drone failures, 
thereby improving mission robustness compared to a single-agent solution \cite{pantelimon2019survey}. 
However, multi-agent systems also introduce challenges related to coordination, task allocation, and mission planning, 
particularly under constraints such as limited communication bandwidth, finite battery capacity, and restricted onboard computation.

\subsection{Coordination in Multi-Agent Systems}

Two common architectural approaches for coordination in multi-agent systems are centralized and decentralized coordination. 
In a centralized architecture, a base station or ground control station is responsible for critical mission decisions and coordinates the behavior of the agents. 
Since the central unit is often equipped with greater computational resources than the individual agents, centralized approaches can support more complex planning and optimization. 
However, this architecture introduces a single point of failure: if the base station fails or communication with it is lost, the mission may be compromised.

In a decentralized architecture, each agent is capable of local decision-making and may communicate directly with other agents. 
This can improve robustness, since the system can remain operational even if individual agents fail. 
Decentralized systems may also scale better as new agents are added, since control is not concentrated in a single central unit. 
However, coordination becomes more challenging because agents must reach decisions using local information and distributed communication mechanisms.

Several mechanisms can support decentralized coordination. 
Leader election methods dynamically select a coordinating agent responsible for organizing group decisions. 
Auction-based methods allocate tasks by allowing agents to bid based on local cost, capability, or availability, with tasks assigned according to the most suitable bid. 
Voting mechanisms allow agents to collectively choose between alternatives, 
while bargaining and negotiation approaches enable agents to iteratively exchange proposals until an agreement is reached.

Hybrid approaches combine elements of both centralized and decentralized coordination. 
For example, a ground station may perform high-level mission planning, 
while individual UAVs retain autonomy to make local decisions based on their current state and surrounding environment. 
Such architectures can exploit the computational resources of centralized systems while preserving some of the robustness and adaptability of decentralized systems.

\subsection{Planning in Multi-Agent Systems}

In multi-agent systems with heterogeneous agents, the task allocation problem consists of assigning tasks to agents under given constraints, such as energy, capability,
and communication limitations, while optimizing a global objective, for example minimizing mission time or maximizing coverage.
Several approaches to multi-agent planning have been identified in the literature \cite{durfee1999dpsp}.

One approach is centralized planning for distributed execution, where a central planner constructs a global plan for all agents. 
This plan specifies both the allocation of tasks and the ordering of actions, and is subsequently distributed to individual agents for execution. 
While this approach simplifies coordination by leveraging a global view of the system, it relies on a central entity and is therefore susceptible to single points of failure.

An alternative approach is distributed planning with centralized plan generation, where multiple agents collaborate to construct a single global plan. 
In this case, agents may contribute specialized knowledge or capabilities during the planning phase, but the resulting plan is still shared among agents for execution. 
This approach can improve flexibility in plan generation, although it still assumes the existence of a coherent global plan.

The most decentralized approach is distributed planning for distributed execution, 
where each agent independently generates its own plan while coordinating with others during execution. 
In such systems, agents may only possess partial knowledge of the overall mission and must dynamically resolve conflicts and dependencies. 
Coordination in this setting often relies on negotiation, communication, or local decision-making mechanisms. 
As a result, a fully global plan may never exist, and system behavior emerges from the interaction of individual agent plans.

In general, centralized planning approaches are simpler to design and can produce globally consistent plans due to their access to complete system information. 
However, decentralized approaches offer greater robustness and scalability, as they avoid reliance on a single coordinating entity. 
The choice between these approaches depends on the requirements of the system, including the need for global optimality, robustness to failures, 
and adaptability to dynamic environments.

Many multi-agent planning problems can be formulated as optimization problems. 
A common abstraction is the multi-traveling salesman problem (mTSP), where multiple agents must visit a set of locations while minimizing a cost function. 
In this formulation, tasks correspond to nodes, agents correspond to salesmen, and decision variables determine the assignment and ordering of tasks for each agent.
While the multi-traveling salesman problem (mTSP) provides a useful abstraction, more realistic formulations are given by the Vehicle Routing Problem (VRP). 
In the VRP, the objective is to determine optimal routes for multiple vehicles that must visit a set of locations. 
In the special case of a single vehicle, the problem reduces to the classical Traveling Salesperson Problem (TSP).
The notion of optimality in VRP depends on the application context. 
A common objective is to minimize the total distance traveled by all vehicles; however, this may lead to unbalanced solutions where a single agent performs most tasks. 
An alternative objective is to minimize the length of the longest route, 
which promotes balanced workload distribution and reduces overall mission completion time. 
Furthermore, VRP formulations can incorporate additional constraints, 
such as capacity limits and time windows, which allow modeling of practical considerations like energy constraints and task deadlines in multi-drone systems.

Different solution approaches exist for such problems. 
Exact methods, such as Mixed Integer Linear Programming (MILP), can provide optimal solutions but often suffer from high computational complexity and limited scalability. 
In contrast, heuristic and metaheuristic methods, such as genetic algorithms, can provide faster solutions and are more suitable for large-scale problems, 
albeit without guarantees of optimality \cite{bektas2006mtsp}.

Finally, real-world multi-agent systems must handle dynamic changes, such as agent failures. 
In such cases, replanning is required, which can be performed either centrally by a global planner or locally by individual agents. 
Hybrid approaches may combine both strategies, allowing systems to balance global optimality with robustness and adaptability.

\section{Large Language Models}
Large Language Models (LLMs) are computational systems designed to generate and understand natural language by modeling the statistical structure of text.
At their core LLMs estimate the probability distribution of the next token given a sequence of preceeding tokens.
Formally, given a context of $\omega_1, \omega_2, \dots, \omega_{t-1}$, the model predicts the next token $w_t$, thereby enabling both language understanding and generation.
The key idea that enables modern LLMs is preptraining, 
where models learn linguistic structure and world knowledge by repeatedly predicting the next token over very large text corpora.

Modern LLMs are implemented using the Transformer architecture introduced in Attention is All You Need \cite{vaswani2017attention}.
The core component of the Transformer is the self-attention mechanism, 
which enables the model to weigh the importance of different tokens relative to each other when generating representations. 
Given an input sequence, attention computes interactions between all token pairs, allowing the model to capture long-range dependencies efficiently. 

Large Language Models vary signficantly in size ranging from millions to hundreds of billions of parameters.
While larger models generally achieve higher performance due to increased representational capacity, they also require substantial computational resources, 
including high memory bandwidth and specialized hardware such as GPUs. 
This poses a challenge for deployment in resource-constrained environments, such as embedded systems, mobile devices, or robotic platforms.
A common approach to address this limitations is quantization, 
where model weights are represented with reduced numerical precision (e.g., 8-bit or 4-bit instead of 32-bit floating point). 
Quantized LLMs significantly reduce memory usage and computational requirements, enabling deployment on edge devices, while maintaining acceptable performance.

\subsection{LLMs in Multi-Agent Systems}
Large Language Models (LLMs) have been successfully applied in various domains, including multi-agent systems, due to their strong reasoning abilities, contextual awareness,
and generalization capabilities.
In this context, their applications can be broadly categorized into communication, perception, planning, and control, following \cite{kim2024llmrobots}.

LLMs are particularly effective at handling the ambiguity of natural language and can be used for both \textit{interpretation} and \textit{grounding}. 
Interpretation refers to converting natural language into structured representations, such as code or formal planning languages (e.g., PDDL). 
Grounding involves mapping human instructions to objects, actions, or processes that can be recognized and executed by agents.

Perception enables agents to operate in real-world environments.
Language models have been extended to multimodal settings by integrating visual and auditory inputs. 
These models, such as Vision-Language Models (VLMs) and Audio Language Models (ALMs), process non-textual data alongside language, 
allowing agents to detect objects and interpret signals from their surroundings.

In planning, traditional approaches rely on predefined algorithms that must be adapted whenever the environment or objectives change. 
In contrast, LLM-based planning through prompting allows for flexible task specification and the incorporation of environmental feedback. 
However, due to the risk of hallucinations, LLM-generated plans should be validated against system constraints or combined with external planning methods.

For control, LLMs can be used in two main ways. 
In the direct approach, the model maps natural language instructions directly to control commands or actions. 
In the indirect approach, the model generates higher-level outputs such as subgoals or plans, which are then executed by separate control systems. 
The indirect approach is generally more flexible and better suited for complex and multi-agent scenarios, 
although it still requires reliable grounding of language into physical actions.


\section{Related Work}

\subsection{LLMs in Autonomous Systems}
Among early LLM-based planners for robotics, SayCan \cite{ahn2022doasican} introduces a framework that grounds language model reasoning in real-world robot capabilities.
In this approach, an LLM generates candidate low-level skills from a high-level natural language instruction.
Each skill is evaluated using two complementary scores: (i) a language model score, which measures how relevant the skill is for completing the task, and (ii) an affordance score, which estimates the feasibility of executing the skill given the robot's current state and environment.
These scores are combined to select the most appropriate action, and the process is iteratively repeated until the task is completed.
This work demonstrates that LLMs can be effectively applied to robotics when their outputs are constrained by physical feasibility.

While SayCan focuses on a single robot, RoCo \cite{mandi2023roco} extends this paradigm to multi-robot systems through collaborative reasoning.
In RoCo, each robot is associated with an LLM-based agent that has access to its own capabilities, local observations, and partial knowledge of the environment.
These agents communicate through a structured dialogue process, iteratively proposing, refining, and critiquing task plans.
Through this interaction, the system converges to a shared global plan, which is then decomposed and distributed among the robots for execution.
Additionally, RoCo leverages LLMs during execution by generating intermediate representations such as 3D waypoints, enabling more flexible and adaptive behavior.

Include research gap in each subsection.
\subsection{Multi-Agent Task Allocation and Coordination}
\subsection{Replanning and Adaptive Systems}